%% 1-7 大きな かずの たし算・ひき算  第5回編集委員会で追加、第11回で 説明を ふやした
%% 第1巻のきまり：15 pt・行送り28 pt、分かち書き（文節の間に半角の空白）、漢字はすべてルビ、式は全角の＋－＝で書く
%% あとの章の けいさん（40－15、50－5、122－119、600－200、500＋500、49＋7、21－18、20－12、90＋90、9＋14）が、
%%   この章の れいだい・じぶんでとこう の形で とけるように した。

%% ---- この章だけで使う絵の部品（頭文字 \bkqv…） ----
\makeatletter
% ひっ算  \bkqvhissan[1マス(mm)]{(x,y)}{けたの数}{記号}{上の数}{下の数}{こたえ}{しるし}
%   (x,y) は 左下のかど。はば（けたの数＋1）×1マス、たかさ 3.6×1マス（いちばん上の 0.6マスは 小さい字の らん）。
%   数は 左の けたから「,」で区切る（なにも かかない ますは 空、書きこみの ますは *）。例 { ,3,6}
%   しるし：けた番号/小さい字/上の数を けすか(1・0)。例 {1/4/1, 2/10/0}。なければ {}。
%   けたの 色は 1-2 の くらいの へやと おなじ（百 みどり、十 あお、一 オレンジ）。
\def\bkqv@star{*}
\newcommand\bkqv@cell[4]{\edef\bkqv@t{#4}%
  \ifx\bkqv@t\bkqv@star
    \filldraw[fill=white, draw=bkOrange, line width=0.9pt, rounded corners=0.8mm] (#2*#1+0.8,#3*#1+0.8) rectangle ++(#1-1.6,#1-1.6);%
  \else
    \node[font=\bkqv@F, text=bkInk, inner sep=0pt] at ({(#2+0.5)*#1},{(#3+0.5)*#1}) {\bkqv@t};%
  \fi}
\newcommand\bkqvhissan[8][9]{\begin{scope}[shift={#2}]
  \pgfmathsetmacro\bkqv@fs{2.5*#1}\pgfmathsetmacro\bkqv@ms{1.35*#1}%
  \def\bkqv@F{\normalfont\fontsize{\bkqv@fs}{\bkqv@fs}\selectfont}%
  \foreach \i in {1,...,#3} {\pgfmathtruncatemacro\bkqv@p{#3-\i}%
    \ifcase\bkqv@p\relax \def\bkqv@bg{bkOrange!13}\or\def\bkqv@bg{bkBlue!11}\else\def\bkqv@bg{bkGreen!12}\fi
    \fill[\bkqv@bg] (\i*#1,0) rectangle ++(#1,3*#1);
    \foreach \j in {0,1,2} {\draw[bkBlue!35, line width=0.5pt] (\i*#1,\j*#1) rectangle ++(#1,#1);}}
  \node[font=\bkqv@F, text=bkInk, inner sep=0pt] at (0.5*#1,1.5*#1) {#4};
  \foreach \d [count=\i] in {#5} {\bkqv@cell{#1}{\i}{2}{\d}}
  \foreach \d [count=\i] in {#6} {\bkqv@cell{#1}{\i}{1}{\d}}
  \foreach \d [count=\i] in {#7} {\bkqv@cell{#1}{\i}{0}{\d}}
  \draw[bkInk, line width=1.4pt, line cap=round] (0.15*#1,#1) -- ({(#3+0.85)*#1},#1);
  \if\relax\detokenize{#8}\relax\else
  \foreach \i/\t/\k in {#8} {
    \node[font=\normalfont\bfseries\fontsize{\bkqv@ms}{\bkqv@ms}\selectfont, text=bkTheme, inner sep=0pt] at ({(\i+0.5)*#1},3.32*#1) {\t};
    \ifnum\k=1 \draw[bkTheme, line width=1.2pt, line cap=round] ({(\i+0.22)*#1},2.18*#1) -- ({(\i+0.78)*#1},2.82*#1);\fi}
  \fi
\end{scope}}
% どんぐり  \bkqvacorn{(x,y)}{たかさ(mm)}   (x,y) は まん中。はばは たかさの 0.7  % noruby
\newcommand\bkqvacorn[2]{\begin{scope}[shift={#1}, scale=#2/10]
  \fill[bkStem!55] (0,-5) .. controls (2.6,-4.7) and (3.4,-1.2) .. (3.2,0.9) -- (-3.2,0.9) .. controls (-3.4,-1.2) and (-2.6,-4.7) .. (0,-5) -- cycle;
  \fill[white, opacity=0.4] (-1.5,-1.5) ellipse (0.55 and 1.4);
  \fill[bkStem] (-3.5,0.5) .. controls (-3.5,4.1) and (3.5,4.1) .. (3.5,0.5) -- cycle;
  \foreach \px in {-2.1,-0.7,0.7,2.1} {\fill[bkStem!70!black] (\px,1.5) circle (0.36);}
  \fill[bkStem!80!black] (-0.4,3.2) rectangle (0.4,4.9);
\end{scope}}
% どんぐり 10この あみぶくろ  \bkqvbag{(x,y)}{どんぐりの たかさ(mm)}   (x,y) は 左下。はば 4.7×、たかさ 3.1×どんぐりの たかさ  % noruby
\newcommand\bkqvbag[2]{\begin{scope}[shift={#1}, scale=#2]
  \fill[bkYellow!14] (0.05,0.25) .. controls (0,0) and (0.1,0) .. (0.35,0) -- (4.35,0) .. controls (4.6,0) and (4.7,0) .. (4.65,0.25) -- (4.55,2.35) -- (2.6,2.7) -- (2.1,2.7) -- (0.15,2.35) -- cycle;
  \foreach \k in {0,...,4} {\bkqvacorn{(0.55+\k*0.9,0.66)}{1} \bkqvacorn{(0.55+\k*0.9,1.72)}{1}}
  \begin{scope}
    \clip (0.05,0.25) .. controls (0,0) and (0.1,0) .. (0.35,0) -- (4.35,0) .. controls (4.6,0) and (4.7,0) .. (4.65,0.25) -- (4.55,2.35) -- (2.6,2.7) -- (2.1,2.7) -- (0.15,2.35) -- cycle;
    \foreach \k in {-3,...,6} {\draw[bkRed!55, line width=0.35pt] (\k*0.75,0) -- (\k*0.75+2.8,2.8) (\k*0.75+2.8,0) -- (\k*0.75,2.8);}
  \end{scope}
  \draw[bkRed!75, line width=0.8pt, line join=round] (0.05,0.25) .. controls (0,0) and (0.1,0) .. (0.35,0) -- (4.35,0) .. controls (4.6,0) and (4.7,0) .. (4.65,0.25) -- (4.55,2.35) -- (2.6,2.7) -- (2.1,2.7) -- (0.15,2.35) -- cycle;
  \fill[bkRed!80, rounded corners=0.5mm] (1.95,2.62) rectangle (2.75,2.95);
  \draw[bkRed!75, line width=0.8pt] (2.2,2.95) -- (2.0,3.15) (2.5,2.95) -- (2.7,3.15);
\end{scope}}
% A・B などの ふだ  \bkqvbadge{(x,y)}{文字}  % noruby
\newcommand\bkqvbadge[2]{\node[fill=bkYellow, text=white, rounded corners=1.5mm, font=\sffamily\bfseries\small, inner xsep=2mm, inner ysep=1mm] at #1 {#2};}
% 数直線（とちゅうの かずから）  \bkqvline{(x,y)}{目もり1つ(mm)}{左はしの かず}{右はしの かず}   (x,y) は 左はしの 目もり  % noruby
\newcommand\bkqvline[4]{\begin{scope}[shift={#1}]
  \pgfmathtruncatemacro\bkqv@n{#4-#3}%
  \draw[bkInk, line width=1pt, dash pattern=on 1.5pt off 1.5pt] (-7,0) -- (-2,0);
  \draw[bkInk, line width=1pt, -{Stealth[length=6pt, width=5pt]}] (-2,0) -- (\bkqv@n*#2+6,0);
  \foreach \i in {0,...,\bkqv@n} {\pgfmathtruncatemacro\bkqv@v{#3+\i}%
    \draw[bkInk, line width=0.8pt] (\i*#2,-1.7) -- (\i*#2,1.7);
    \node[anchor=north, font=\sffamily\bfseries\footnotesize, text=bkInk, inner sep=0pt] at (\i*#2,-2.8) {\bkqv@v};}
\end{scope}}
% 数直線の 上を すすむ やじるし（1つずつ 番号つき）  \bkqvhops{(x,y)}{目もり1つ(mm)}{いくつ}{色}   (x,y) は はじめの 目もり  % noruby
\newcommand\bkqvhops[4]{\begin{scope}[shift={#1}]
  \foreach \i in {1,...,#3} {\draw[#4, line width=1.3pt, -{Stealth[length=5pt]}] ({(\i-1)*#2+0.7},1.2) .. controls ({(\i-0.8)*#2},7.5) and ({(\i-0.2)*#2},7.5) .. ({\i*#2-0.7},1.4);
    \node[circle, fill=white, draw=#4, line width=0.6pt, font=\sffamily\bfseries\scriptsize, text=#4, inner sep=0.5pt, minimum size=3.6mm] at ({(\i-0.5)*#2},8.6) {\i};}
  \fill[#4] (0,0) circle (1.5);
\end{scope}}
% きょりメーター  \bkqvodo{(x,y)}{4つの すうじ}{この ばんごうから 右の わを きいろに（0なら なし）}   はば 38、たかさ 13  % noruby
\newcommand\bkqvodo[3]{\begin{scope}[shift={#1}]
  \fill[bkInk!85, rounded corners=1.5mm] (0,0) rectangle (38,13);
  \foreach \d [count=\k] in {#2} {\pgfmathtruncatemacro\bkqv@hl{(#3>0)&&(\k>=#3)}%
    \ifnum\bkqv@hl=1 \fill[bkYellow!60] (\k*8.6-6.2,2) rectangle ++(7.2,9);\else\fill[white] (\k*8.6-6.2,2) rectangle ++(7.2,9);\fi
    \node[font=\sffamily\bfseries\fontsize{16}{16}\selectfont, text=bkInk] at (\k*8.6-2.6,6.5) {\d};}
\end{scope}}
\makeatother

\bkchapter{1}{7}{\R{大}{おお}きな かずの\\たし\R{算}{ざん}・ひき\R{算}{ざん}}{\bklevel{1}}{\R{大}{おお}きな かずの たし\R{算}{ざん}・ひき\R{算}{ざん}}
\bkmichishirube{1-2,1-4,1-5,1-6}{1-10,1-21,2-3}

\begin{goal}
\begin{itemize}[label={\color{bkTheme}●}, leftmargin=1.4\zw, itemsep=0pt, topsep=0pt]
\item \R{何十}{なん|じゅう}・\R{何百}{なん|びゃく}の けいさんが できる。
\item くらいを そろえて、 ひっ\R{算}{さん}が できる。
\item くり\R{上}{あ}がり・くり\R{下}{さ}がりの わけが わかる。
\end{itemize}
\end{goal}

\begin{hatena}
\begin{bkfigure}
\begin{tikzpicture}[x=1mm, y=1mm]
  % 1くみ：ふくろ 3つと ばら 6こ  % noruby
  \fill[white, rounded corners=3mm] (0,0) rectangle (59,45);
  \bkqvbadge{(10,39)}{1くみ}
  \node[anchor=east, font=\sffamily\bfseries, text=bkStem!90!black] at (55,39) {36こ};
  \foreach \x in {2,20.8,39.6} {\bkqvbag{(\x,17)}{3.7}}
  \foreach \x in {14,20.3,26.6,32.9,39.2,45.5} {\bkqvacorn{(\x,8)}{6}}
  % 2くみ：ふくろ 2つと ばら 7こ  % noruby
  \begin{scope}[shift={(63,0)}]
  \fill[white, rounded corners=3mm] (0,0) rectangle (59,45);
  \bkqvbadge{(10,39)}{2くみ}
  \node[anchor=east, font=\sffamily\bfseries, text=bkStem!90!black] at (55,39) {27こ};
  \foreach \x in {11.1,30.5} {\bkqvbag{(\x,17)}{3.7}}
  \foreach \x in {10.6,16.9,23.2,29.5,35.8,42.1,48.4} {\bkqvacorn{(\x,8)}{6}}
  \end{scope}
\end{tikzpicture}
\end{bkfigure}
\noindent どんぐりを 10こずつ ふくろに いれたよ。\\
1くみと 2くみ、 あわせて なんこ？ ゆびでは たりない！\\
ふくろは いくつ、 ばらは いくつに なるかな？
\end{hatena}

\clearpage
%% ---------------- p.77 やってみよう ----------------
\begin{yattemiyou}[10の ぼうと ばらで 36 ＋ 27]
\begin{enumerate}[label={\bkcircled[bkGreen]{\arabic*}}, leftmargin=2.2\zw, itemsep=2pt, topsep=0pt]
\item 10の ぼう（1-2）と ばらで、 36と 27を つくろう。
\end{enumerate}
\begin{bkfigure}
\begin{tikzpicture}[x=1mm, y=1mm]
  \fill[bkBlue!12, rounded corners=2mm] (18,0) rectangle (66,57);
  \fill[bkOrange!14, rounded corners=2mm] (68,0) rectangle (120,57);
  \node[font=\sffamily\bfseries\small, text=bkBlue!80!black] at (42,52.5) {\R{十}{じゅう}の くらい};
  \node[font=\sffamily\bfseries\small, text=bkOrange!85!black] at (94,52.5) {\R{一}{いち}の くらい};
  \draw[bkGray!70, line width=0.7pt, dash pattern=on 2pt off 1.5pt] (2,24.5) -- (120,24.5);
  % 36
  \node[font=\normalfont\fontsize{22}{22}\selectfont, text=bkInk] at (9,37) {36};
  \foreach \k in {0,1,2} {\bkrod{(33.4+\k*6,26)}{2.2}{bkRoute4}}
  \foreach \k in {0,...,5} {\bkcube{(77.2+\k*5.6,35.9)}{2.2}{bkRoute4}}
  % 27
  \node[font=\normalfont\fontsize{22}{22}\selectfont, text=bkInk] at (9,12.5) {27};
  \foreach \k in {0,1} {\bkrod{(36.4+\k*6,1.5)}{2.2}{bkRoute4}}
  \foreach \k in {0,...,6} {\bkcube{(74.4+\k*5.6,11.4)}{2.2}{bkRoute4}}
\end{tikzpicture}
\end{bkfigure}
\begin{enumerate}[label={\bkcircled[bkGreen]{\arabic*}}, leftmargin=2.2\zw, itemsep=2pt, topsep=0pt, start=2]
\item ぼうは ぼう、 ばらは ばらで あつめよう。\par
ぼうは\enspace\bkkakikomi\enspace \R{本}{ほん}、\hspace{1\zw}ばらは\enspace\bkkakikomi\enspace こ
\item ばらが 10こより おおいね。 ばら 10こを、 10の ぼう 1\R{本}{ぽん}と とりかえよう。
\end{enumerate}
\begin{bkfigure}
\begin{tikzpicture}[x=1mm, y=1mm]
  \foreach \i in {0,...,9} {\pgfmathtruncatemacro\cx{mod(\i,5)}\pgfmathtruncatemacro\cy{div(\i,5)}
    \bkcube{(\cx*3.6+2,\cy*3.6+8)}{2}{bkRoute4}}
  \foreach \i in {0,1,2} {\bkcube{(\i*3.6+25,8)}{2}{bkRoute4}}
  \draw[bkOrange, line width=1pt, dash pattern=on 2pt off 1.5pt, rounded corners=2mm] (0,6) rectangle (20.4,15.6);
  \node[anchor=north, font=\sffamily\bfseries\footnotesize, text=bkInk] at (17.2,4) {ばら 13こ};
  \draw[bk arrow, bkOrange, line width=1.2pt] (38,11) -- (54,11);
  \node[font=\sffamily\bfseries\footnotesize, text=bkOrange!85!black] at (46,15.5) {とりかえる};
  \bkrod{(62,1)}{2}{bkRoute4}
  \draw[bkOrange, line width=1pt, dash pattern=on 2pt off 1.5pt, rounded corners=1.5mm] (60,-0.5) rectangle (66,22.5);
  \foreach \i in {0,1,2} {\bkcube{(\i*3.6+71,8)}{2}{bkRoute4}}
  \node[anchor=west, font=\sffamily\bfseries\footnotesize, text=bkInk] at (83,9) {ぼう 1\R{本}{ぽん}と ばら 3こ};
\end{tikzpicture}
\end{bkfigure}
\begin{enumerate}[label={\bkcircled[bkGreen]{\arabic*}}, leftmargin=2.2\zw, itemsep=2pt, topsep=0pt, start=4]
\item ぼうは\enspace\bkkakikomi\enspace \R{本}{ほん}、\hspace{1\zw}ばらは\enspace\bkkakikomi\enspace こ。\par
ぜんぶで\enspace\bkkakikomi\enspace こ
\end{enumerate}
\end{yattemiyou}

\bkfukidashi{mina}{ばら 10こが、 ぼう 1\R{本}{ぽん}に かわったね。 ブロックが なければ、 えを ゆびで さして かぞえても いいよ。}

\clearpage
%% ---------------- 数学の道具（第11回：ブロックの図と 説明の段を ふやした） ----------------
\bksection{doguu}{\R{数学}{すう|がく}の \R{道具}{どう|ぐ}}

\bksubsection{1}{10・100の まとまりで けいさん}

\label{bk:1-7:matomari}30 ＋ 40は いくつ？ 30は 10の ぼうが 3\R{本}{ぼん}、 40は 4\R{本}{ほん}です（1-2）。

\begin{bkfigure}
\begin{tikzpicture}[x=1mm, y=1mm]
  \foreach \k in {0,1,2} {\bkrod{(\k*5,5)}{2}{bkRoute4}}
  \node[font=\fontsize{20}{20}\selectfont\bfseries, text=bkTheme] at (19,15) {＋};
  \foreach \k in {0,...,3} {\bkrod{(26+\k*5,5)}{2}{bkRoute2}}
  \node[font=\fontsize{20}{20}\selectfont\bfseries, text=bkTheme] at (50,15) {→};
  \foreach \k in {0,1,2} {\bkrod{(57+\k*5,5)}{2}{bkRoute4}}
  \foreach \k in {3,...,6} {\bkrod{(57+\k*5,5)}{2}{bkRoute2}}
  \node[anchor=north, font=\sffamily\bfseries\footnotesize, text=bkRoute4!70!black] at (6,3.5) {10が 3こ};
  \node[anchor=north, font=\sffamily\bfseries\footnotesize, text=bkRoute2!80!black] at (35.5,3.5) {10が 4こ};
  \node[anchor=north, font=\sffamily\bfseries\footnotesize, text=bkInk] at (73,3.5) {10が 7こ};
  \node[anchor=west, font=\sffamily\bfseries, text=bkTheme] at (94,15) {70};
\end{tikzpicture}
\end{bkfigure}
\noindent ぼうの \R{本数}{ほん|すう}だけを たせば よいので、 3 ＋ 4 ＝ 7。 10が 7こで 70です。
\bkbigshiki{30 ＋ 40 ＝ 70}{10が 3 ＋ 4 ＝ 7こ}

\noindent{\sffamily\bfseries\color{bkTheme}10が 10こ より おおく なったら}\par
90 ＋ 90は、 10が 9 ＋ 9 ＝ 18こ。 10の ぼうが 10\R{本}{ぽん} あつまると、 100の いた 1まいと とりかえられます（1-2）。

\begin{bkfigure}
\begin{tikzpicture}[x=1mm, y=1mm]
  \foreach \k in {0,...,8} {\bkrod{(\k*3,0)}{1.8}{bkRoute4}}
  \foreach \k in {9,...,17} {\bkrod{(\k*3,0)}{1.8}{bkRoute2}}
  \draw[bkOrange, line width=1pt, dash pattern=on 2pt off 1.5pt, rounded corners=1.5mm] (-1.2,-1.2) rectangle (29.4,19.2);
  \node[anchor=north, font=\sffamily\bfseries\footnotesize, text=bkInk] at (26.4,-2.4) {10が 18こ};
  \draw[bk arrow, bkOrange, line width=1.2pt] (55,9) -- (65,9);
  \node[font=\sffamily\bfseries\footnotesize, text=bkOrange!85!black] at (60,23) {とりかえる};
  \bkflat{(68,0)}{1.8}{bkRoute6}
  \draw[bkOrange, line width=1pt, dash pattern=on 2pt off 1.5pt, rounded corners=1.5mm] (66.8,-1.2) rectangle (87.2,19.2);
  \foreach \k in {0,...,7} {\bkrod{(90+\k*3,0)}{1.8}{bkRoute2}}
  \node[anchor=north, font=\sffamily\bfseries\footnotesize, text=bkInk] at (94,-2.4) {100と 80で 180};
\end{tikzpicture}
\end{bkfigure}
\bkbigshiki{90 ＋ 90 ＝ 180}{10が 18こ}

\clearpage
\noindent{\sffamily\bfseries\color{bkTheme}100の まとまりでも おなじ}\par
600 － 200は、 100の いたで かんがえます。 100が 6まい から 2まい とります。

\begin{bkfigure}
\begin{tikzpicture}[x=1mm, y=1mm]
  \foreach \k in {0,...,5} {\bkflat{(\k*16,0)}{1.3}{bkRoute6}}
  \foreach \k in {4,5} {\draw[bkRed, line width=1.6pt] (\k*16-1,-1) -- (\k*16+14,14) (\k*16-1,14) -- (\k*16+14,-1);}
  \bkbrace{(0,-1.5)}{(61,-1.5)}{のこりは 100が 4まい}{-1}
  \node[anchor=west, font=\sffamily\bfseries\footnotesize, text=bkRed!85!black] at (99,6.5) {2まい とる};
\end{tikzpicture}
\end{bkfigure}
\noindent 100が 6 － 2 ＝ 4まいで、 400です。 \mbox{500 ＋ 500}なら、 100が 5 ＋ 5 ＝ 10まい。 100が 10まいで 1000です（1-2）。\bkidx{なんじゅう}{\R{何十}{なん|じゅう}・\R{何百}{なん|びゃく}の けいさん}

\bksubsection{2}{くらいごとに たす・ひく}

\label{bk:1-7:kuraigoto}24 ＋ 13は、 10の ぼうは ぼうどうし、 ばらは ばらどうしで たそう。

\begin{bkfigure}
\begin{tikzpicture}[x=1mm, y=1mm]
  \node[font=\normalfont\fontsize{18}{18}\selectfont, text=bkRoute4!75!black] at (6,22) {24};
  \node[font=\normalfont\fontsize{18}{18}\selectfont, text=bkRoute2!85!black] at (6,12) {13};
  \fill[bkBlue!12, rounded corners=2mm] (16,0) rectangle (64,33);
  \fill[bkOrange!14, rounded corners=2mm] (66,0) rectangle (120,33);
  \node[font=\sffamily\bfseries\small, text=bkBlue!80!black] at (40,29.5) {\R{十}{じゅう}の くらい};
  \node[font=\sffamily\bfseries\small, text=bkOrange!85!black] at (93,29.5) {\R{一}{いち}の くらい};
  \foreach \k in {0,1} {\bkrod{(31+\k*5,7.5)}{1.8}{bkRoute4}}
  \bkrod{(45,7.5)}{1.8}{bkRoute2}
  \foreach \k in {0,...,3} {\bkcube{(73+\k*5,15)}{2}{bkRoute4}}
  \foreach \k in {0,1,2} {\bkcube{(95+\k*5,15)}{2}{bkRoute2}}
  \node[font=\sffamily\bfseries\small, text=bkInk] at (40,3.5) {2 ＋ 1 ＝ 3};
  \node[font=\sffamily\bfseries\small, text=bkInk] at (93,3.5) {4 ＋ 3 ＝ 7};
\end{tikzpicture}
\end{bkfigure}
\noindent ぼうと ばらを まぜて たすと、 「10」と 「1」を おなじに かぞえて しまいます。 だから、 おなじ くらいどうしで たします。
\bkbigshiki{24 ＋ 13 ＝ 37}{10が 3こと 1が 7こ}

\clearpage
\noindent{\sffamily\bfseries\color{bkTheme}ひき\R{算}{ざん}も くらいごとに}\par
38 － 15は、 38の ブロックから、 ぼう 1\R{本}{ぽん}と ばら 5こを とります。

\begin{bkfigure}
\begin{tikzpicture}[x=1mm, y=1mm]
  \fill[bkBlue!12, rounded corners=2mm] (16,0) rectangle (64,29.5);
  \fill[bkOrange!14, rounded corners=2mm] (66,0) rectangle (120,29.5);
  \node[font=\normalfont\fontsize{18}{18}\selectfont, text=bkInk] at (6,16) {38};
  \node[font=\sffamily\bfseries\small, text=bkBlue!80!black] at (40,26.3) {\R{十}{じゅう}の くらい};
  \node[font=\sffamily\bfseries\small, text=bkOrange!85!black] at (93,26.3) {\R{一}{いち}の くらい};
  \foreach \k in {0,1,2} {\bkrod{(31+\k*6,5)}{1.8}{bkRoute4}}
  \draw[bkRed, line width=1.4pt] (41.5,4) -- (47.5,24) (47.5,4) -- (41.5,24);
  \foreach \k in {0,...,7} {\bkcube{(71+\k*6,14)}{2.4}{bkRoute4}}
  \foreach \k in {3,...,7} {\draw[bkRed, line width=1.3pt] (70.4+\k*6,13.4) -- (74+\k*6,17) (70.4+\k*6,17) -- (74+\k*6,13.4);}
  \node[font=\sffamily\bfseries\small, text=bkInk] at (40,2.8) {3 － 1 ＝ 2};
  \node[font=\sffamily\bfseries\small, text=bkInk] at (93,2.8) {8 － 5 ＝ 3};
\end{tikzpicture}
\end{bkfigure}
\noindent のこりは、 ぼう 2\R{本}{ほん}と ばら 3こ。 10が 2こと 1が 3こで 23です。

\bksubsection{3}{ひっ\R{算}{さん}と くり\R{上}{あ}がり}

\label{bk:1-7:kuriagari}\bkhatenaref の 36 ＋ 27を、 ブロックの かわりに すうじで けいさんしよう。

\noindent\begin{minipage}[c]{0.34\linewidth}\centering
\begin{tikzpicture}[x=1mm, y=1mm]
  \bkqvhissan[9]{(0,0)}{2}{＋}{3,6}{2,7}{6,3}{1/1/0}
  \node[anchor=south, font=\sffamily\bfseries\footnotesize, text=bkTheme, inner sep=1pt] at (9,35.5) {ぼう 1\R{本}{ぽん}};
  \draw[bkTheme, line width=0.8pt, -{Stealth[length=3.5pt]}] (11.5,35.6) -- (13,32.6);
\end{tikzpicture}
\end{minipage}\hfill
\begin{minipage}[c]{0.64\linewidth}
\begin{enumerate}[label={\bkcircled[bkTheme]{\arabic*}}, leftmargin=2\zw, itemsep=0pt, topsep=0pt]
\item くらいを たてに そろえて かく。
\item \R{一}{いち}の くらい。 6 ＋ 7 ＝ 13（1-4）
\item 13は 10と 3。 3を かき、 10は ぼう 1\R{本}{ぽん}に して \R{十}{じゅう}の くらいへ。 \R{小}{ちい}さく 1と かく。
\item \R{十}{じゅう}の くらい。 1 ＋ 3 ＋ 2 ＝ 6
\end{enumerate}
\end{minipage}
\par\vspace{2pt}
\noindent この かきかたを ひっ\R{算}{さん}と いいます。\bkidx{ひっさん}{ひっ\R{算}{さん}} 10を 1つ \R{左}{ひだり}の くらいへ うつす ことを、 くり\R{上}{あ}がりと いいます。\bkidx{くりあがり}{くり\R{上}{あ}がり} \R{小}{ちい}さい 1は 「ぼうが 1\R{本}{ぽん} ふえたよ」の メモです。 \bkhatenaref の こたえは 63こです。

\clearpage
\bksubsection{4}{ひっ\R{算}{さん}と くり\R{下}{さ}がり}

\label{bk:1-7:kurisagari}52 － 27を、 ブロックで \R{見}{み}て みよう。

\begin{bkfigure}
\begin{tikzpicture}[x=1mm, y=1mm]
  % A：52
  \fill[bkPaper, rounded corners=2mm] (0,0) rectangle (36,31);
  \node[circle, fill=bkTheme, text=white, font=\sffamily\bfseries\scriptsize, inner sep=0pt, minimum size=4.6mm] at (4.5,27) {1};
  \foreach \k in {0,...,4} {\bkrod{(4+\k*3.4,5)}{1.8}{bkRoute4}}
  \foreach \k in {0,1} {\bkcube{(24+\k*3,5)}{2.2}{bkRoute4}}
  \node[font=\sffamily\bfseries\footnotesize] at (18,1.8) {52};
  \draw[bk arrow, bkGray, line width=1pt] (37.5,15) -- (44,15);
  % B：ぼう1本を ばら10こに
  \fill[bkPaper, rounded corners=2mm] (45,0) rectangle (85,31);
  \node[circle, fill=bkTheme, text=white, font=\sffamily\bfseries\scriptsize, inner sep=0pt, minimum size=4.6mm] at (49.5,27) {2};
  \foreach \k in {0,...,3} {\bkrod{(49+\k*3.4,5)}{1.8}{bkRoute4}}
  \foreach \k in {0,...,11} {\pgfmathtruncatemacro\cx{mod(\k,4)}\pgfmathtruncatemacro\cy{div(\k,4)}
    \ifnum\k<10 \bkcube{(65+\cx*3,5+\cy*3)}{2.2}{bkOrange}\else\bkcube{(65+\cx*3,5+\cy*3)}{2.2}{bkRoute4}\fi}
  \node[font=\sffamily\bfseries\footnotesize] at (65,1.8) {40と 12};
  \draw[bk arrow, bkGray, line width=1pt] (86.5,15) -- (93,15);
  % C：27を とる
  \fill[bkPaper, rounded corners=2mm] (94,0) rectangle (134,31);
  \node[circle, fill=bkTheme, text=white, font=\sffamily\bfseries\scriptsize, inner sep=0pt, minimum size=4.6mm] at (98.5,27) {3};
  \foreach \k in {0,...,3} {\bkrod{(98+\k*3.4,5)}{1.8}{bkRoute4}}
  \foreach \k in {2,3} {\draw[bkRed, line width=1.2pt] (97.4+\k*3.4,4) -- (100.6+\k*3.4,24) (100.6+\k*3.4,4) -- (97.4+\k*3.4,24);}
  \foreach \k in {0,...,11} {\pgfmathtruncatemacro\cx{mod(\k,4)}\pgfmathtruncatemacro\cy{div(\k,4)}
    \ifnum\k<10 \bkcube{(114+\cx*3,5+\cy*3)}{2.2}{bkOrange}\else\bkcube{(114+\cx*3,5+\cy*3)}{2.2}{bkRoute4}\fi
    \ifnum\k<7 \draw[bkRed, line width=1pt] (113.8+\cx*3,4.8+\cy*3) -- (116.4+\cx*3,7.4+\cy*3);\fi}
  \node[font=\sffamily\bfseries\footnotesize] at (114,1.8) {のこり 25};
\end{tikzpicture}
\end{bkfigure}
\begin{enumerate}[label={\bkcircled[bkGreen]{\arabic*}}, leftmargin=2.2\zw, itemsep=0pt, topsep=0pt]
\item 52は、 ぼう 5\R{本}{ほん}と ばら 2こ。 ばらから 7こは とれない。
\item ぼう 1\R{本}{ぽん}を ばら 10こに くずす（1-5）。 ぼうは 4\R{本}{ほん}、 ばらは 10と 2で 12こ。
\item ばらから 7こ、 ぼうから 2\R{本}{ほん} とる。 のこりは ぼう 2\R{本}{ほん}と ばら 5こで、 25。
\end{enumerate}
\noindent{\sffamily\bfseries\color{bkTheme}ひっ\R{算}{さん}では}\quad おなじ ことを、 すうじで こう かきます。

\noindent\begin{minipage}[c]{0.34\linewidth}\centering
\begin{tikzpicture}[x=1mm, y=1mm]
  \bkqvhissan[9]{(0,0)}{2}{－}{5,2}{2,7}{2,5}{1/4/1, 2/10/0}
\end{tikzpicture}
\end{minipage}\hfill
\begin{minipage}[c]{0.64\linewidth}
\begin{enumerate}[label={\bkcircled[bkTheme]{\arabic*}}, leftmargin=2\zw, itemsep=0pt, topsep=0pt]
\item \R{一}{いち}の くらい。 2から 7は とれない。
\item \R{十}{じゅう}の くらいから 10を もらう。 5を けして 4（ぼうが 1\R{本}{ぽん} へった）、 \R{一}{いち}の くらいの \R{上}{うえ}に 10。
\item 10と 2で 12。 12 － 7 ＝ 5（1-5）
\item \R{十}{じゅう}の くらい。 4 － 2 ＝ 2
\end{enumerate}
\end{minipage}
\par\vspace{2pt}
\noindent 1つ \R{左}{ひだり}の くらいから 10を もらう ことを、 くり\R{下}{さ}がりと いいます。\bkidx{くりさがり}{くり\R{下}{さ}がり} こたえは、 25 ＋ 27 ＝ 52で たしかめよう（1-5）。

\clearpage
\bksubsection{5}{100を こえる たし\R{算}{ざん}}

\label{bk:1-7:hyaku}86 ＋ 57は、 \R{十}{じゅう}の くらいでも くり\R{上}{あ}がりが おきます。

\noindent\begin{minipage}[c]{0.3\linewidth}\centering
\begin{tikzpicture}[x=1mm, y=1mm]
  \bkqvhissan[8.5]{(0,0)}{3}{＋}{ ,8,6}{ ,5,7}{1,4,3}{1/1/0, 2/1/0}
\end{tikzpicture}
\end{minipage}\hfill
\begin{minipage}[c]{0.68\linewidth}
\begin{enumerate}[label={\bkcircled[bkTheme]{\arabic*}}, leftmargin=2\zw, itemsep=0pt, topsep=0pt]
\item \R{一}{いち}の くらい。 6 ＋ 7 ＝ 13。 3を かき、 1を \R{十}{じゅう}の くらいへ くり\R{上}{あ}げる。
\item \R{十}{じゅう}の くらい。 1 ＋ 8 ＋ 5 ＝ 14。 これは 10の ぼうが 14\R{本}{ほん}。
\item ぼう 10\R{本}{ぽん}は 100の いた 1まい。 4を かき、 1を \R{百}{ひゃく}の くらいへ くり\R{上}{あ}げる。
\item \R{百}{ひゃく}の くらいは 1。 こたえは 143。
\end{enumerate}
\end{minipage}

\begin{bkfigure}
\begin{tikzpicture}[x=1mm, y=1mm]
  \foreach \k in {0,...,13} {\bkrod{(\k*3,0)}{1.8}{bkRoute2}}
  \draw[bkOrange, line width=1pt, dash pattern=on 2pt off 1.5pt, rounded corners=1.5mm] (-1.2,-1.2) rectangle (29.4,19.2);
  \node[anchor=north, font=\sffamily\bfseries\footnotesize, text=bkInk] at (20.4,-2.4) {10が 14こ};
  \draw[bk arrow, bkOrange, line width=1.2pt] (46,9) -- (57,9);
  \node[font=\sffamily\bfseries\footnotesize, text=bkOrange!85!black] at (51.5,13.5) {とりかえる};
  \bkflat{(61,0)}{1.8}{bkRoute6}
  \foreach \k in {0,...,3} {\bkrod{(84+\k*3,0)}{1.8}{bkRoute2}}
  \node[anchor=north, font=\sffamily\bfseries\footnotesize, text=bkInk] at (80,-2.4) {100が 1こと 10が 4こ};
\end{tikzpicture}
\end{bkfigure}
\noindent どの くらいでも、 10 あつまったら 1つ \R{左}{ひだり}の くらいへ 1を くり\R{上}{あ}げます。 \R{一}{いち}の くらいでも、 \R{十}{じゅう}の くらいでも、 おなじ きまりです。

\begin{taisetsu}
くらいを そろえて、 \R{一}{いち}の くらいから けいさん。\\
10に なったら、 \R{左}{ひだり}の くらいへ くり\R{上}{あ}がり。\\
とれない ときは、 \R{左}{ひだり}から 10を もらう（くり\R{下}{さ}がり）。
\end{taisetsu}

\clearpage
\bksubsection{6}{こたえを たしかめよう}

\label{bk:1-7:tashikame}ひっ\R{算}{さん}の こたえが あって いるかは、 ぎゃくの けいさんで たしかめられます。

\noindent{\sffamily\bfseries\color{bkTheme}ひき\R{算}{ざん}は、 たし\R{算}{ざん}で たしかめる}\par
52 － 27 ＝ 25 でした。 52を、 とった 27と のこりの 25に \R{分}{わ}けた テープで \R{見}{み}て みよう。

\begin{bkfigure}
\begin{tikzpicture}[x=1mm, y=1mm]
  \bktape{(0,0)}{50}{bkBlue}{のこり 25}
  \bktape{(50,0)}{54}{bkOrange}{とった 27}
  \bkbrace{(0,8)}{(104,8)}{はじめの 52}{1}
\end{tikzpicture}
\end{bkfigure}
\noindent のこりの 25と、 とった 27を あわせれば、 はじめの 52に もどる はずです（1-6の 「もどって たしかめよう」と おなじ）。
\bkbigshiki{25 ＋ 27 ＝ 52}{52に もどった\enspace\bkmaru}

\noindent{\sffamily\bfseries\color{bkTheme}たし\R{算}{ざん}は、 ひき\R{算}{ざん}か、 \R{入}{い}れかえで たしかめる}\par
36 ＋ 27 ＝ 63 なら、 63 － 27 ＝ 36で もとに もどります。 27 ＋ 36の ように、 たす じゅんばんを \R{入}{い}れかえても、 こたえは おなじ 63に なります。

\noindent{\sffamily\bfseries\color{bkTheme}やって みよう}\par
\noindent\hspace*{1\zw}\begin{tblr}{colspec={Q[l,m,wd=74mm] Q[l,m,wd=52mm]}, colsep=0pt, rowsep=4pt}
40 － 15 ＝ 25 の たしかめ & 25 ＋ 15 ＝\enspace\bkkakikomi[16mm] \\
86 ＋ 57 ＝ 143 の たしかめ & 143 － 57 ＝\enspace\bkkakikomi[16mm] \\
\end{tblr}
\par\vspace{2pt}
\noindent もとの 40や 86に もどれば、 こたえは あって います。

\clearpage
%% ---------------- 例題 ----------------

\begin{reidai}[\R{何十}{なん|じゅう}・\R{何百}{なん|びゃく}の けいさん]
（1）50 ＋ 80\hspace{4\zw}（2）600 － 200

\kangaekata （1）10が 5 ＋ 8 ＝ 13こ。 10が 10こで 100だから、 13こは 100と 30。\\
（2）100が 6 － 2 ＝ 4こ。

\kotae （1）130\hspace{3\zw}（2）400
\end{reidai}

\begin{jibun}
\noindent\begin{tblr}{colspec={Q[l,m,wd=64mm] Q[l,m,wd=64mm]}, colsep=0pt, rowsep=3pt}
（1）90 ＋ 90 ＝\enspace\bkkakikomi[18mm] & （2）500 ＋ 500 ＝\enspace\bkkakikomi[18mm] \\
（3）700 － 400 ＝\enspace\bkkakikomi[18mm] & \hspace{0.6\zw}{\small\color{bkOrange!58!black}ヒント：100が 10こで 1000} \\
\end{tblr}
\bkjans{（1）180\quad（2）1000\quad（3）300}
\end{jibun}

\begin{reidai}[100を こえる たし\R{算}{ざん}]
\noindent\begin{minipage}[c]{0.76\linewidth}
ビーズが 75こと 48こ。 あわせて なんこ？\par\vspace{4pt}

\kangaekata \R{一}{いち}の くらい 5 ＋ 8 ＝ 13。 3を かいて 1 くり\R{上}{あ}げる。\\
\R{十}{じゅう}の くらい 1 ＋ 7 ＋ 4 ＝ 12。 2を かいて、 1を \R{百}{ひゃく}の くらいへ。
\end{minipage}\hfill
\begin{minipage}[c]{0.22\linewidth}\raggedleft
\begin{tikzpicture}[x=1mm, y=1mm]
  \bkqvhissan[7]{(0,0)}{3}{＋}{ ,7,5}{ ,4,8}{1,2,3}{1/1/0, 2/1/0}
\end{tikzpicture}
\end{minipage}
\par\vspace{5pt}
\shiki 75 ＋ 48 ＝ 123 \kotaenext 123こ
\end{reidai}

\begin{jibun}[ひっ\R{算}{さん}で しよう]
\begin{bkfigure}
\begin{tikzpicture}[x=1mm, y=1mm]
  \foreach \px/\l in {0/（1）,42/（2）,84/（3）} {\node[anchor=north west, font=\small, text=bkInk, inner sep=0pt] at (\px,22) {\l};}
  \bkqvhissan[6.5]{(10,0)}{3}{＋}{ ,3,8}{ ,4,6}{*,*,*}{}
  \bkqvhissan[6.5]{(52,0)}{3}{＋}{ ,6,7}{ ,5,8}{*,*,*}{}
  \bkqvhissan[6.5]{(94,0)}{3}{＋}{ ,5,9}{ ,6,4}{*,*,*}{}
\end{tikzpicture}
\end{bkfigure}
\bkjans{（1）84\quad（2）125\quad（3）123}
\end{jibun}

\begin{reidai}[くり\R{下}{さ}がりの ある ひき\R{算}{ざん}]
いろがみが 40まい。 15まい つかうと、 のこりは？\par\vspace{3pt}

\noindent\begin{minipage}[c]{0.8\linewidth}
\kangaekata \R{一}{いち}の くらいは 0。 0から 5は とれない。\\
\R{十}{じゅう}の くらいから 10を もらう。 4は 3に なる。\\
10 － 5 ＝ 5。 \R{十}{じゅう}の くらいは 3 － 1 ＝ 2。

\shiki 40 － 15 ＝ 25 \kotaenext 25まい
\end{minipage}\hfill
\begin{minipage}[c]{0.19\linewidth}\raggedleft
\begin{tikzpicture}[x=1mm, y=1mm]
  \bkqvhissan[7]{(0,0)}{2}{－}{4,0}{1,5}{2,5}{1/3/1, 2/10/0}
\end{tikzpicture}
\end{minipage}
\begin{bkfigure}
\begin{tikzpicture}[x=1mm, y=1mm]
  \foreach \k in {0,...,3} {\bkrod{(\k*3.4,0)}{1.6}{bkRoute4}}
  \node[anchor=north, font=\sffamily\bfseries\scriptsize] at (6,-1) {40};
  \draw[bk arrow, bkGray, line width=1pt] (17,8) -- (25,8);
  \node[font=\sffamily\bfseries\scriptsize, text=bkGray] at (21,12) {くずす};
  \foreach \k in {0,...,2} {\bkrod{(29+\k*3.4,0)}{1.6}{bkRoute4}}
  \foreach \k in {0,...,9} {\pgfmathtruncatemacro\cx{mod(\k,5)}\pgfmathtruncatemacro\cy{div(\k,5)}\bkcube{(41+\cx*2.8,\cy*2.8)}{2}{bkOrange}}
  \node[anchor=north, font=\sffamily\bfseries\scriptsize] at (42,-1) {30と 10};
  \draw[bk arrow, bkGray, line width=1pt] (58,8) -- (66,8);
  \node[font=\sffamily\bfseries\scriptsize, text=bkGray] at (62,12) {15 とる};
  \foreach \k in {0,...,2} {\bkrod{(70+\k*3.4,0)}{1.6}{bkRoute4}}
  \draw[bkRed, line width=1.1pt] (69.5,-0.5) -- (72.1,16.5) (72.1,-0.5) -- (69.5,16.5);
  \foreach \k in {0,...,9} {\pgfmathtruncatemacro\cx{mod(\k,5)}\pgfmathtruncatemacro\cy{div(\k,5)}\bkcube{(82+\cx*2.8,\cy*2.8)}{2}{bkOrange}
    \ifnum\k<5 \draw[bkRed, line width=0.9pt] (81.8+\cx*2.8,-0.2+\cy*2.8) -- (84.2+\cx*2.8,2.2+\cy*2.8);\fi}
  \node[anchor=west, font=\sffamily\bfseries\footnotesize] at (98,4) {のこり 25};
\end{tikzpicture}
\end{bkfigure}
\end{reidai}

\begin{jibun}[ひっ\R{算}{さん}で しよう]
\begin{bkfigure}
\begin{tikzpicture}[x=1mm, y=1mm]
  \foreach \px/\l in {0/（1）,42/（2）,84/（3）} {\node[anchor=north west, font=\small, text=bkInk, inner sep=0pt] at (\px,22) {\l};}
  \bkqvhissan[6.5]{(14,0)}{2}{－}{6,3}{2,8}{*,*}{}
  \bkqvhissan[6.5]{(56,0)}{2}{－}{7,0}{3,6}{*,*}{}
  \bkqvhissan[6.5]{(98,0)}{2}{－}{8,1}{2,4}{*,*}{}
\end{tikzpicture}
\end{bkfigure}
\bkjans{（1）35\quad（2）34\quad（3）57}
\end{jibun}

\begin{reidai}[けたの かずが ちがう ひっ\R{算}{さん}]
（1）49 ＋ 7\hspace{4\zw}（2）50 － 5\par\vspace{3pt}

\noindent\begin{minipage}[c]{0.66\linewidth}
\kangaekata 7も 5も \R{一}{いち}の くらいの かず。 \R{右}{みぎ}はしを そろえて かく。 \R{十}{じゅう}の くらいの \R{下}{した}は、 なにも ない（0と おなじ。 1-3）。\\
（1）9 ＋ 7 ＝ 16。 1を くり\R{上}{あ}げて、 1 ＋ 4 ＝ 5。\\
（2）0から 5は とれない。 10を もらって 10 － 5 ＝ 5。 \R{十}{じゅう}の くらいは 4。
\end{minipage}\hfill
\begin{minipage}[c]{0.32\linewidth}\raggedleft
\begin{tikzpicture}[x=1mm, y=1mm]
  \bkqvhissan[6.5]{(0,0)}{2}{＋}{4,9}{ ,7}{5,6}{1/1/0}
  \bkqvhissan[6.5]{(23,0)}{2}{－}{5,0}{ ,5}{4,5}{1/4/1, 2/10/0}
\end{tikzpicture}
\end{minipage}
\par\vspace{3pt}
\kotae （1）56\hspace{3\zw}（2）45
\end{reidai}

\begin{jibun}[ひっ\R{算}{さん}で しよう]
\begin{bkfigure}
\begin{tikzpicture}[x=1mm, y=1mm]
  \foreach \px/\l in {0/（1）,33/（2）,66/（3）,99/（4）} {\node[anchor=north west, font=\small, text=bkInk, inner sep=0pt] at (\px,22) {\l};}
  \bkqvhissan[6.5]{(8,0)}{2}{＋}{ ,9}{1,4}{*,*}{}
  \bkqvhissan[6.5]{(41,0)}{2}{＋}{3,6}{ ,8}{*,*}{}
  \bkqvhissan[6.5]{(74,0)}{2}{－}{6,0}{ ,4}{*,*}{}
  \bkqvhissan[6.5]{(107,0)}{2}{－}{4,3}{ ,6}{*,*}{}
\end{tikzpicture}
\end{bkfigure}
\bkjans{（1）23\quad（2）44\quad（3）56\quad（4）37}
\end{jibun}

\begin{reidai}[ちがいが \R{小}{ちい}さい ひき\R{算}{ざん}]
ソラは 101かい、 ミナは 96かい とんだ。 ちがいは？
\begin{bkfigure}
\begin{tikzpicture}[x=1mm, y=1mm]
  \bkqvline{(8,0)}{12}{94}{102}
  \bkqvhops{(32,0)}{12}{5}{bkBlue}
  \node[anchor=east, font=\sffamily\bfseries\footnotesize, text=bkAppleRed!80!black] at (29,5) {ミナ};
  \node[anchor=west, font=\sffamily\bfseries\footnotesize, text=bkRoute2!80!black] at (95,5) {ソラ};
  \fill[bkRoute2] (92,0) circle (1.5);
\end{tikzpicture}
\end{bkfigure}
\noindent\begin{minipage}[c]{0.72\linewidth}
\kangaekata ひっ\R{算}{さん}に すると（\R{右}{みぎ}）、 くり\R{下}{さ}がりが 2かいも あって たいへん。 2つの かずが \R{近}{ちか}い ときは、 96から 101まで かぞえたす（1-6の ちがい）。 97、 98、 99、 100、 101で 5かい。
\end{minipage}\hfill
\begin{minipage}[c]{0.26\linewidth}\raggedleft
\begin{tikzpicture}[x=1mm, y=1mm]
  \bkqvhissan[7]{(0,0)}{3}{－}{1,0,1}{ ,9,6}{ , ,5}{1/0/1, 2/9/1, 3/10/0}
\end{tikzpicture}
\end{minipage}

\shiki 101 － 96 ＝ 5 \kotaenext 5かい
\end{reidai}

\begin{jibun}[\R{本}{ほん}を よんだ ページ]
ミナは 103ページ、 ソラは 97ページ。 ちがいは？ \R{数直線}{すう|ちょく|せん}に とんだ やじるしを かいて かぞえよう。
\begin{bkfigure}
\begin{tikzpicture}[x=1mm, y=1mm]
  \bkqvline{(8,0)}{12}{96}{104}
\end{tikzpicture}
\end{bkfigure}
\bkshikiran[9\zw]
\bkjans{103 － 97 ＝ 6（97から 103まで かぞえたす）\quad 6ページ}
\end{jibun}

\begin{reidai}[\R{百}{ひゃく}の くらいから くり\R{下}{さ}がる ひき\R{算}{ざん}]
\noindent\begin{minipage}[c]{0.74\linewidth}
135 － 52を ひっ\R{算}{さん}で しよう。\par\vspace{3pt}

\kangaekata \R{一}{いち}の くらい 5 － 2 ＝ 3。\\
\R{十}{じゅう}の くらいは、 3から 5は とれない。 \R{百}{ひゃく}の くらいから 100（10の ぼう 10\R{本}{ぽん}）を もらう。 10と 3で 13。 13 － 5 ＝ 8。\\
\R{百}{ひゃく}の くらいは 1を けして 0。 なにも かかない。
\end{minipage}\hfill
\begin{minipage}[c]{0.24\linewidth}\raggedleft
\begin{tikzpicture}[x=1mm, y=1mm]
  \bkqvhissan[7]{(0,0)}{3}{－}{1,3,5}{ ,5,2}{ ,8,3}{1/0/1, 2/10/0}
\end{tikzpicture}
\end{minipage}
\par\vspace{3pt}
\begin{bkfigure}
\begin{tikzpicture}[x=1mm, y=1mm]
  \bkflat{(0,0)}{1.3}{bkRoute6}
  \foreach \k in {0,1,2} {\bkrod{(16+\k*3,0)}{1.3}{bkRoute4}}
  \foreach \k in {0,...,4} {\bkcube{(26+\k*2.4,0)}{1.8}{bkOrange}}
  \node[anchor=north, font=\sffamily\bfseries\scriptsize] at (19,-1) {135};
  \draw[bk arrow, bkGray, line width=1pt] (40,6.5) -- (49,6.5);
  \node[anchor=east, font=\sffamily\bfseries\scriptsize, text=bkGray] at (51,11) {100を くずす};
  \foreach \k in {0,...,12} {\ifnum\k<10 \bkrod{(53+\k*2.8,0)}{1.3}{bkRoute6}\else\bkrod{(53+\k*2.8,0)}{1.3}{bkRoute4}\fi
    \ifnum\k>7 \draw[bkRed, line width=1pt] (52.6+\k*2.8,-0.4) -- (54.7+\k*2.8,13.4);\fi}
  \foreach \k in {0,...,4} {\bkcube{(91+\k*2.4,0)}{1.8}{bkOrange}
    \ifnum\k>2 \draw[bkRed, line width=0.9pt] (90.8+\k*2.4,-0.2) -- (93+\k*2.4,2);\fi}
  \node[anchor=north, font=\sffamily\bfseries\scriptsize] at (78,-1) {10が 13こと 5。 52を とる};
  \node[anchor=west, font=\sffamily\bfseries\footnotesize] at (105,6.5) {のこり 83};
\end{tikzpicture}
\end{bkfigure}
\shiki 135 － 52 ＝ 83 \kotaenext 83\qquad{\small\color{bkGray}たしかめ：83 ＋ 52 ＝ 135}
\end{reidai}

\begin{jibun}[ひっ\R{算}{さん}で しよう]
\begin{bkfigure}
\begin{tikzpicture}[x=1mm, y=1mm]
  \foreach \px/\l in {0/（1）,42/（2）,84/（3）} {\node[anchor=north west, font=\small, text=bkInk, inner sep=0pt] at (\px,22) {\l};}
  \bkqvhissan[6.5]{(10,0)}{3}{－}{1,2,8}{ ,7,4}{*,*,*}{}
  \bkqvhissan[6.5]{(52,0)}{3}{－}{1,1,6}{ ,3,5}{*,*,*}{}
  \bkqvhissan[6.5]{(94,0)}{3}{－}{1,5,3}{ ,9,2}{*,*,*}{}
\end{tikzpicture}
\end{bkfigure}
\bkjans{（1）54\quad（2）81\quad（3）61}
\end{jibun}

\begin{reidai}[たし\R{算}{ざん}かな？ ひき\R{算}{ざん}かな？]
（1）\R{赤}{あか}い \R{花}{はな}が 38\R{本}{ぽん}、 \R{白}{しろ}い \R{花}{はな}が 25\R{本}{ほん}。 あわせて なん\R{本}{ぼん}？\\
（2）64ページの \R{本}{ほん}を 27ページ よんだ。 のこりは なんページ？\\
（3）どんぐりを ミナは 73こ、 ソラは 48こ ひろった。 ちがいは なんこ？
\begin{bkfigure}
\begin{tikzpicture}[x=1mm, y=1mm]
  \node[anchor=west, font=\sffamily\bfseries\footnotesize] at (0,13) {（1）あわせる};
  \bktape{(0,0)}{19}{bkAppleRed}{38}
  \bktape{(19,0)}{13}{bkBlue}{25}
  \bkbrace{(0,-1)}{(32,-1)}{？}{-1}
  \begin{scope}[shift={(44,0)}]
  \node[anchor=west, font=\sffamily\bfseries\footnotesize] at (0,13) {（2）のこり};
  \bktape{(0,0)}{32}{bkGreen}{}
  \fill[bkGray!35] (0,0.4) rectangle (13.5,6.6);
  \node[font=\sffamily\bfseries\scriptsize] at (6.75,3.5) {27};
  \node[font=\sffamily\bfseries\scriptsize] at (22.75,3.5) {？};
  \bkbrace{(0,-1)}{(32,-1)}{64}{-1}
  \end{scope}
  \begin{scope}[shift={(88,0)}]
  \node[anchor=west, font=\sffamily\bfseries\footnotesize] at (0,13) {（3）ちがい};
  \bktape{(0,3.5)}{36.5}{bkOrange}{73}
  \bktape{(0,-4)}{24}{bkRoute6}{48}
  \draw[bkRed, line width=1.1pt, {Stealth[length=3pt]}-{Stealth[length=3pt]}] (24.5,-0.5) -- (36,-0.5);
  \node[anchor=north, font=\sffamily\bfseries\scriptsize, text=bkRed] at (30.2,-5) {？};
  \end{scope}
\end{tikzpicture}
\end{bkfigure}
\kangaekata 「あわせて」は たし\R{算}{ざん}（1-4）、 「のこり」「ちがい」は ひき\R{算}{ざん}（1-5）。

\shiki （1）38 ＋ 25 ＝ 63\quad（2）64 － 27 ＝ 37\quad（3）73 － 48 ＝ 25

\kotae （1）63\R{本}{ぼん}\quad（2）37ページ\quad（3）25こ
\end{reidai}

\begin{jibun}
（1）バスに 34\R{人}{にん}。 18\R{人}{にん} のって きた。 みんなで なん\R{人}{にん}？
\bkshikiran[9\zw]
（2）おりがみ 50まいから 23まい つかった。 のこりは？
\bkshikiran[9\zw]
\bkjans{（1）34 ＋ 18 ＝ 52\quad 52\R{人}{にん}\quad（2）50 － 23 ＝ 27\quad 27まい}
\end{jibun}


\clearpage
%% ---------------- p.82 まちがえやすいところ ----------------
\bksection{tsumazuki}{まちがえやすい ところ}

\begin{tsumazuki}{くらいが ずれる}
\noindent\begin{minipage}[c]{0.44\linewidth}
\begin{tikzpicture}[x=1mm, y=1mm]
  \node at (2,10) {\bkbatsu};
  \bkqvhissan[6]{(5,0)}{2}{＋}{3,6}{2, }{5,6}{}
  \node at (30,10) {\bkmaru};
  \bkqvhissan[6]{(33,0)}{2}{＋}{3,6}{ ,2}{3,8}{}
\end{tikzpicture}
\end{minipage}\hfill
\begin{minipage}[c]{0.54\linewidth}
36 ＋ 2の 2は、 \R{一}{いち}の くらい。 \R{右}{みぎ}はしを そろえて かこう。
\end{minipage}
\end{tsumazuki}

\begin{tsumazuki}{くり\R{上}{あ}がりの 1を わすれる}
\noindent\begin{minipage}[c]{0.44\linewidth}
\begin{tikzpicture}[x=1mm, y=1mm]
  \node at (2,10) {\bkbatsu};
  \bkqvhissan[6]{(5,0)}{2}{＋}{3,6}{2,7}{5,3}{}
  \node at (30,10) {\bkmaru};
  \bkqvhissan[6]{(33,0)}{2}{＋}{3,6}{2,7}{6,3}{1/1/0}
\end{tikzpicture}
\end{minipage}\hfill
\begin{minipage}[c]{0.54\linewidth}
\R{小}{ちい}さく かいた 1も、 \R{十}{じゅう}の くらいで たそう。
\end{minipage}
\end{tsumazuki}

\begin{tsumazuki}{\R{下}{した}の かずから ひいて しまう}
\noindent\begin{minipage}[c]{0.44\linewidth}
\begin{tikzpicture}[x=1mm, y=1mm]
  \node at (2,10) {\bkbatsu};
  \bkqvhissan[6]{(5,0)}{2}{－}{5,2}{2,7}{3,5}{}
  \node at (30,10) {\bkmaru};
  \bkqvhissan[6]{(33,0)}{2}{－}{5,2}{2,7}{2,5}{1/4/1, 2/10/0}
\end{tikzpicture}
\end{minipage}\hfill
\begin{minipage}[c]{0.54\linewidth}
7 － 2を して しまった。 \R{上}{うえ}の 2から 7は とれないので、 10を もらおう。
\end{minipage}
\end{tsumazuki}

\begin{column}{どうして \R{一}{いち}の くらいから？}
36 ＋ 27を \R{十}{じゅう}の くらいから すると、 3 ＋ 2 ＝ 5で、 さきに 5と かきます。 でも \R{一}{いち}の くらいで くり\R{上}{あ}がるので、 5を 6に かきなおす ことに なります。 \R{一}{いち}の くらいから すれば、 かきなおさずに すみます。
\end{column}

\clearpage
%% ---------------- p.83 つかわれているところ ----------------
\begin{butsuri}
\begin{bkfigure}
\begin{tikzpicture}[x=1mm, y=1mm]
  % きょりメーター
  \fill[white, rounded corners=3mm] (0,0) rectangle (59,44);
  \node[font=\sffamily\bfseries\small] at (29.5,39.5) {きょりメーター};
  \bkqvodo{(10.5,22)}{0,1,9,9}{0}
  \bkqvodo{(10.5,3)}{0,2,0,0}{2}
  \draw[bk arrow, bkTheme] (29.5,20.5) -- (29.5,16.5);
  \node[anchor=west, font=\sffamily\bfseries\fontsize{8.5}{9}\selectfont, text=bkTheme] at (31,18.5) {1 すすむと};
  % コンピューター
  \begin{scope}[shift={(63,0)}]
  \fill[white, rounded corners=3mm] (0,0) rectangle (59,44);
  \node[font=\sffamily\bfseries\small] at (29.5,39.5) {コンピューターの \R{中}{なか}};
  \node[font=\sffamily\bfseries\footnotesize, text=bkInk] at (29.5,32.6) {86 ＋ 57 ＝};
  \foreach \x in {10,14,...,50} {\fill[bkGray!70] (\x-0.8,2) rectangle (\x+0.8,5) (\x-0.8,26) rectangle (\x+0.8,29);}
  \fill[bkInk!80, rounded corners=1.5mm] (6,5) rectangle (53,26);
  \foreach \x/\d in {10/1,24/4,38/3} {\fill[white, rounded corners=1mm] (\x,7) rectangle ++(11,10);
    \node[font=\sffamily\bfseries\fontsize{15}{15}\selectfont, text=bkInk] at (\x+5.5,12) {\d};}
  \foreach \x in {24,38} {\draw[bkYellow, line width=1.2pt, -{Stealth[length=4pt]}] (\x+6,18) .. controls (\x+3,23.5) and (\x-5,23.5) .. (\x-7.5,18);
    \node[font=\sffamily\bfseries\fontsize{8}{9}\selectfont, text=bkYellow] at (\x-1,24.2) {1};}
  \end{scope}
\end{tikzpicture}
\end{bkfigure}
\begin{itemize}[label={\color{bkPurple}●}, leftmargin=1.4\zw, itemsep=0pt, topsep=0pt]
\item \R{車}{くるま}の きょりメーターの わは、 9から 0に なる ときだけ、 となりの わを つめで 1つ おします。 0199の つぎは 0200です。
\item でんたくや コンピューターも、 くらいごとに たして、 くり\R{上}{あ}がりを となりへ わたします。 コンピューターは 0と 1だけで かくので、 2で くり\R{上}{あ}がります（2\R{巻}{かん}）。
\end{itemize}
\end{butsuri}

\bkfukidashi{delta}{くり\R{上}{あ}がりが あるから、 どんなに \R{大}{おお}きな かずでも たせるんだ。 「たしあわせの \R{路線}{ろ|せん}」を すすもう！}

\begin{bkfigure}
\begin{tikzpicture}[x=1mm, y=1mm]
  \draw[bkRoute2, line width=3pt] (10,0) -- (128,0);
  \node[anchor=west, inner sep=0pt] at (0,12) {\bkrouteicon[1.4em]{2}\hspace{0.3em}{\sffamily\bfseries\small\color{bkRoute2!80!black}たしあわせの \R{路線}{ろ|せん}}};
  \foreach \x/\c/\t/\s in {10/bkVol1/{1-4}/{たし\R{算}{ざん}},39.5/bkVol1/{1-7}/{\R{大}{おお}きな かず},69/bkVol1/{1-12}/{かけ\R{算}{ざん}},98.5/bkVol1/{1-21}/{ひろさ},128/bkVol2/{2-3}/{{\ltjsetparameter{jacharrange={+200}}＋－×÷}}} {
    \fill[white, draw=\c, line width=1.6pt] (\x,0) circle (2.6);
    \node[anchor=base, font=\sffamily\bfseries\footnotesize, text=\c!85!black] at (\x,-8.4) {\t};
    \node[anchor=base, font=\sffamily\footnotesize, text=bkInk] at (\x,-15.4) {\s};}
  \node[fill=bkTheme, text=white, rounded corners=1mm, font=\sffamily\bfseries\fontsize{8}{9}\selectfont, inner sep=1.5pt] at (39.5,5.5) {いま ここ};
\end{tikzpicture}
\end{bkfigure}

\clearpage
%% ---------------- p.84〜85 もんだい ----------------
\bkendproblems

\bkproblemgroup{きほん}{1}

\bkproblem けいさんを しましょう。
\par\vspace{2pt}
\noindent\hspace*{1\zw}\begin{tblr}{colspec={*{2}{Q[l,m,wd=66mm]}}, colsep=0pt, rowsep=3pt}
（1）40 ＋ 30 ＝ \bkkakikomi[16mm] & （2）70 ＋ 60 ＝ \bkkakikomi[16mm] \\
（3）90 － 50 ＝ \bkkakikomi[16mm] & （4）300 ＋ 400 ＝ \bkkakikomi[16mm] \\
（5）600 ＋ 400 ＝ \bkkakikomi[16mm] & （6）800 － 500 ＝ \bkkakikomi[16mm] \\
\end{tblr}
\bkpans{（1）70\quad（2）130\quad（3）40\quad（4）700\quad（5）1000\quad（6）300}

\bkproblem ひっ\R{算}{さん}で しましょう。
\begin{bkfigure}
\begin{tikzpicture}[x=1mm, y=1mm]
  \foreach \px/\l in {0/（1）,44/（2）,88/（3）} {\node[anchor=north west, font=\small, text=bkInk, inner sep=0pt] at (\px,26) {\l};}
  \bkqvhissan[7.5]{(9,0)}{3}{＋}{ ,2,5}{ ,3,4}{*,*,*}{}
  \bkqvhissan[7.5]{(53,0)}{3}{＋}{ ,4,7}{ ,3,6}{*,*,*}{}
  \bkqvhissan[7.5]{(97,0)}{3}{＋}{ , ,8}{ ,6,5}{*,*,*}{}
\end{tikzpicture}
\end{bkfigure}
\bkpans{（1）59\quad（2）83\quad（3）73}

\bkproblem ひっ\R{算}{さん}で しましょう。
\begin{bkfigure}
\begin{tikzpicture}[x=1mm, y=1mm]
  \foreach \px/\l in {0/（1）,44/（2）,88/（3）} {\node[anchor=north west, font=\small, text=bkInk, inner sep=0pt] at (\px,26) {\l};}
  \bkqvhissan[7.5]{(13,0)}{2}{－}{6,8}{2,5}{*,*}{}
  \bkqvhissan[7.5]{(57,0)}{2}{－}{7,4}{3,8}{*,*}{}
  \bkqvhissan[7.5]{(101,0)}{2}{－}{6,0}{2,3}{*,*}{}
\end{tikzpicture}
\end{bkfigure}
\bkpans{（1）43\quad（2）36\quad（3）37}

\bkproblem \begin{minipage}[t]{0.66\linewidth}
どんぐりを、 ソラは 37こ、 ミナは 45こ ひろいました。 あわせて なんこ？
\bkshikiran[4.5\zw]
\end{minipage}\hfill
\begin{minipage}[t]{0.24\linewidth}\raggedleft
\begin{tikzpicture}[x=1mm, y=1mm, baseline={([yshift=-3mm]current bounding box.north)}]
  \node[anchor=south west, font=\sffamily\bfseries\scriptsize, text=bkOrange!85!black, inner sep=0pt] at (0,25.6) {ひっ\R{算}{さん}};
  \bkqvhissan[7]{(0,0)}{3}{}{ , , }{ , , }{ , , }{}
\end{tikzpicture}
\end{minipage}
\bkpans{37 ＋ 45 ＝ 82\quad 82こ}

\bkproblem \begin{minipage}[t]{0.66\linewidth}
シールが 52まい あります。 18まい つかうと、 のこりは なんまい？
\bkshikiran[4.5\zw]
\end{minipage}\hfill
\begin{minipage}[t]{0.24\linewidth}\raggedleft
\begin{tikzpicture}[x=1mm, y=1mm, baseline={([yshift=-3mm]current bounding box.north)}]
  \node[anchor=south west, font=\sffamily\bfseries\scriptsize, text=bkOrange!85!black, inner sep=0pt] at (0,25.6) {ひっ\R{算}{さん}};
  \bkqvhissan[7]{(0,0)}{3}{}{ , , }{ , , }{ , , }{}
\end{tikzpicture}
\end{minipage}
\bkpans{52 － 18 ＝ 34\quad 34まい}

\bkproblemgroup{チャレンジ}{1}

\bkproblem \begin{minipage}[t]{0.66\linewidth}
なわとびを、 きのう 68かい、 きょう 57かい とびました。 あわせて なんかい？
\bkshikiran[4.5\zw]
\end{minipage}\hfill
\begin{minipage}[t]{0.24\linewidth}\raggedleft
\begin{tikzpicture}[x=1mm, y=1mm, baseline={([yshift=-3mm]current bounding box.north)}]
  \node[anchor=south west, font=\sffamily\bfseries\scriptsize, text=bkOrange!85!black, inner sep=0pt] at (0,25.6) {ひっ\R{算}{さん}};
  \bkqvhissan[7]{(0,0)}{3}{}{ , , }{ , , }{ , , }{}
\end{tikzpicture}
\end{minipage}
\bkpans{68 ＋ 57 ＝ 125\quad 125かい}

\bkproblem ミナの クラスは 102さつ、 ソラの クラスは 98さつ \R{本}{ほん}を よみました。 ちがいは なんさつ？
\bkshikiran
\bkpans{102 － 98 ＝ 4\quad 4さつ}

\begin{furikaeri}
\begin{checklist}
\item \R{何十}{なん|じゅう}・\R{何百}{なん|びゃく}の けいさんが できる。\hfill{\footnotesize\sffamily\color{bkTheme}→\pageref{bk:1-7:matomari}ページ}
\item くらいを そろえて、 ひっ\R{算}{さん}が できる。\hfill{\footnotesize\sffamily\color{bkTheme}→\pageref{bk:1-7:kuraigoto}ページ}
\item くり\R{上}{あ}がり・くり\R{下}{さ}がりが わかる。\hfill{\footnotesize\sffamily\color{bkTheme}→\pageref{bk:1-7:kuriagari}ページ}
\end{checklist}
チェックが つかなかったら、 →の ページに もどろう。
\end{furikaeri}

\bkfukidashi{sora}{かいものの レシートを \R{見}{み}せて もらって、 ねだんを ひっ\R{算}{さん}で たして みよう！}
\clearpage
